\section{Case Study: Multimodal AI Integration at Rang Dong}

The e-commerce system of Rang Dong Light Source and Vacuum Flask Joint Stock Company, Vietnam’s oldest lighting manufacturer, has integrated Multimodal AI to power its advanced Sales Agent. By processing a catalog of over 1,500 product variants with high technical complexity—such as IoT-enabled smart home systems and configurable LED solutions—the system uses multimodal capabilities to bridge the gap between customer intent and technical specifications.

\subsection{Multimodal Application: Vision and Language}

The system’s core strength lies in its Image-based Identification combined with natural language processing. Customers interacting through the e-commerce interface can upload photographs of products encountered in showrooms or construction sites. The AI identifies visually similar items from the catalog and allows users to apply simultaneous text-based constraints, such as finding a specific product but with "higher power". This dual-modality approach resolves the limitations of visual matching alone, where two identical-looking lights may have vastly different industrial certifications or internal components.

\subsection{Operational Improvements and Results}

The application of this technology has yielded significant improvements in digital commerce efficiency:
\begin{itemize}
    \item \textbf{Reduced Search Friction:} The system achieved a \textbf{95.61\%} quality score in product search and \textbf{99.76\%} in model code lookup, allowing users to find precise technical documentation in seconds.
    
    \item \textbf{Adoption and Engagement:} Over five months, the system processed \textbf{5,797} requests across \textbf{3,121} unique sessions on the official website.
    
    \item \textbf{Conversational Depth:} Approximately \textbf{35.6\%} of sessions involved multi-turn consultations, indicating that users found the automated advisory useful enough to sustain interactions.
    
    \item \textbf{Service Scalability:} The AI absorbed an estimated \textbf{260 to 416 hours} of human-equivalent labor.
    
    \item \textbf{Temporal Extension:} Notably, \textbf{46\%} of interactions occurred outside standard business hours, capturing demand that would otherwise require expensive 24/7 staffing.
\end{itemize}